\chapter{Konsequenzen: Einheit der Physik}

Die WDBT+ ist keine alternative Theorie, sondern die \emph{fundamentalere Beschreibung} der physikalischen Realität.

\subsection{Die Hierarchie der Theorien}

\begin{itemize}
    \item Die \textbf{Quantenmechanik} ist der Spezialfall \( D = 3 \) der fraktalen Wellengleichung.
    \item Die \textbf{Quantenelektrodynamik} ist die effektive Theorie, die durch die fraktale Regularisierung der WDBT+ endlich wird.
    \item Die \textbf{Allgemeine Relativitätstheorie} ist der Grenzfall \( D \to 3 \) der \( \Omega \)-Theorie, die Gravitation als Rotation beschreibt.
\end{itemize}

\subsection{Falsifizierbarkeit}

Die Theorie macht testbare Vorhersagen:

\begin{enumerate}
    \item Die fraktale Gewichtung der Atomorbitale führt zu einer charakteristischen Aufspaltung der Energieniveaus, die von der Spin-Bahn-Kopplung unterscheidbar ist.
    \item Der Lamb-Shift weicht um den Faktor \( (a_0/l_*)^{D-3} \) von der QED-Vorhersage ab.
    \item Die effektive Hubble-Konstante ist skalenabhängig: \( H_{\text{eff}}(d) \propto d^{0,704} \).
\end{enumerate}

\subsection{Die Einheit der Physik}

Die fraktale Dimension \( D \approx 2,704 \) ist die fundamentale Konstante, die alle Bereiche der Physik miteinander verbindet – von der Quantenmechanik über die QED bis zur Kosmologie.